
\documentclass{article}
\usepackage{colortbl}
\usepackage{makecell}
\usepackage{multirow}
\usepackage{supertabular}

\begin{document}

\newcounter{utterance}

\centering \large Interaction Transcript for game `cladder', experiment `full\_v1.5\_default', episode 8797 with selfrefine{-}llama{-}3.1{-}8b{-}instruct.
\vspace{24pt}

{ \footnotesize  \setcounter{utterance}{1}
\setlength{\tabcolsep}{0pt}
\begin{supertabular}{c@{$\;$}|p{.15\linewidth}@{}p{.15\linewidth}p{.15\linewidth}p{.15\linewidth}p{.15\linewidth}p{.15\linewidth}}
    \# & \multicolumn{2}{c}{Player} && \multicolumn{2}{c}{Game Master} \\
    \hline

    \theutterance \stepcounter{utterance}  
    & & & \multicolumn{4}{p{0.6\linewidth}}{
        \cellcolor[rgb]{0.9,0.9,0.9}{
            \makecell[{{p{\linewidth}}}]{
                \texttt{\tiny{[P1$\langle$GM]}}
                \texttt{You are an expert in causal inference. The following question is not a typical commonsense query, but rather a meticulously designed question created by a professor specializing in causal inference, intended to assess the students' mastery of the course content.} \\
\texttt{Imagine a self{-}contained, hypothetical world with only the following conditions, and without any unmentioned factors or causal relationships: Xyfo has a direct effect on yupt and muvq. Yupt has a direct effect on zupj. Zupj has a direct effect on muvq. Xyfo is unobserved. For those who are not yupt and are not zupj, the probability of muvq is 90\%. For those who are not yupt and are zupj, the probability of muvq is 54\%. For those who are yupt and are not zupj, the probability of muvq is 96\%. For those who are yupt and are zupj, the probability of muvq is 57\%. For those who are not yupt, the probability of zupj is 51\%. For those who are yupt, the probability of zupj is 34\%. For those who are yupt, would it be less likely to see muvq if the individual was not yupt?} \\
\texttt{Start your answer with 'Yes' or 'No', followed by additional reasoning or evidence to support your explanation.} \\
            }
        }
    }
     \\ \\

    \theutterance \stepcounter{utterance}  
    & & \multicolumn{4}{p{0.6\linewidth}}{
        \cellcolor[rgb]{0.95,0.95,0.95}{
            \makecell[{{p{\linewidth}}}]{
                \texttt{\tiny{[P1$\leftrightarrow$P]}}
                \texttt{{[}DRAFT{]}} \\
\texttt{No.} \\
\\ 
\texttt{To assess the causal effect of yupt on muvq, we can use the backdoor adjustment formula, which is given by:} \\
\\ 
\texttt{P(muvq | yupt) = P(muvq | yupt, zupj) * P(zupj | yupt) + P(muvq | yupt, not zupj) * P(not zupj | yupt)} \\
\\ 
\texttt{We are given the following probabilities:} \\
\\ 
\texttt{{-} P(muvq | yupt, zupj) = 57\%} \\
\texttt{{-} P(muvq | yupt, not zupj) = 96\%} \\
\texttt{{-} P(zupj | yupt) = 34\%} \\
\texttt{{-} P(not zupj | yupt) = 66\%} \\
\\ 
\texttt{Substituting these values into the formula, we get:} \\
\\ 
\texttt{P(muvq | yupt) = 0.57 * 0.34 + 0.96 * 0.66} \\
\texttt{P(muvq | yupt) = 0.194 + 0.6336} \\
\texttt{P(muvq | yupt) = 0.8276} \\
\\ 
\texttt{Now, we need to compare this with the probability of muvq for those who are not yupt. We are given that for those who are not yupt, the probability of muvq is 90\%. However, we also need to consider the effect of zupj on muvq for those who are not yupt. We are given that for those who are not yupt and are zupj, the probability of muvq is 54\%, and for those who are not yupt and are not zupj, the probability of muvq is 90\%.} \\
\\ 
\texttt{To calculate the overall probability of muvq for those who are not yupt, we need to use the law of total probability:} \\
\\ 
\texttt{P(muvq | not yupt) = P(muvq | not yupt, zupj) * P(zupj | not yupt) + P(muvq | not yupt, not zupj) * P(not zupj | not yupt)} \\
\\ 
\texttt{We are given that for those who are not yupt, the probability of zupj is 51\%. Therefore, P(not zupj | not yupt) = 49\%. Substituting the given values, we get:} \\
\\ 
\texttt{P(muvq | not yupt) = 0.54 * 0.51 + 0.90 * 0.49} \\
\texttt{P(muvq | not yupt) = 0.2754 + 0.441} \\
\texttt{P(muvq | not yupt) = 0.7164} \\
\\ 
\texttt{Comparing the probabilities, we see that P(muvq | yupt) = 0.8276 is greater than P(muvq | not yupt) = 0.7164. This suggests that it is more likely to see} \\
            }
        }
    }
    & \\ \\

    \theutterance \stepcounter{utterance}  
    & & \multicolumn{4}{p{0.6\linewidth}}{
        \cellcolor[rgb]{0.95,0.95,0.95}{
            \makecell[{{p{\linewidth}}}]{
                \texttt{\tiny{[P1$\leftrightarrow$P]}}
                \texttt{{[}SELF{-}FEEDBACK{]}} \\
\texttt{FORMAT: ok} \\
\texttt{RULES: ok} \\
\texttt{STRATEGY: problem} \\
\texttt{REASON: The answer incorrectly concludes that it is more likely to see muvq if the individual was not yupt, when in fact the calculated probabilities show the opposite.} \\
            }
        }
    }
    & \\ \\

    \theutterance \stepcounter{utterance}  
    & \multicolumn{4}{p{0.6\linewidth}}{
        \cellcolor[rgb]{0.561,0.737,0.561}{
            \makecell[{{p{\linewidth}}}]{
                \texttt{\tiny{[P1$\rangle$GM]}}
                \texttt{No.} \\
\\ 
\texttt{The calculated probabilities show that P(muvq | yupt) = 0.8276 is greater than P(muvq | not yupt) = 0.7164. This suggests that it is more likely to see muvq if the individual was yupt.} \\
            }
        }
    }
    & & \\ \\

    \theutterance \stepcounter{utterance}  
    & & & \multicolumn{2}{p{0.3\linewidth}}{
        \cellcolor[rgb]{0.9,0.9,0.9}{
            \makecell[{{p{\linewidth}}}]{
                \texttt{\tiny{[GM$|$GM]}}
                \texttt{No} \\
            }
        }
    }
    & & \\ \\

    \theutterance \stepcounter{utterance}  
    & & & \multicolumn{2}{p{0.3\linewidth}}{
        \cellcolor[rgb]{0.9,0.9,0.9}{
            \makecell[{{p{\linewidth}}}]{
                \texttt{\tiny{[GM$|$GM]}}
                \texttt{yes} \\
            }
        }
    }
    & & \\ \\

    \theutterance \stepcounter{utterance}  
    & & & \multicolumn{2}{p{0.3\linewidth}}{
        \cellcolor[rgb]{0.9,0.9,0.9}{
            \makecell[{{p{\linewidth}}}]{
                \texttt{\tiny{[GM$|$GM]}}
                \texttt{game\_result = LOSE} \\
            }
        }
    }
    & & \\ \\

\end{supertabular}
}

\end{document}
